\documentclass[11pt,a4paper]{article}

% Paquetes
\usepackage[utf8]{inputenc}
\usepackage[T1]{fontenc}
\usepackage[spanish]{babel}
\usepackage{geometry}
\usepackage{hyperref}
\usepackage{enumitem}
\usepackage{titlesec}
\usepackage{fancyhdr}
\usepackage{lastpage}

% Fuente serif clásica
\usepackage{mathpazo} % Palatino

% Configuración de página
\geometry{left=2.5cm, right=2.5cm, top=2.5cm, bottom=2.5cm}

% Hipervínculos discretos
\hypersetup{
    colorlinks=true,
    linkcolor=black,
    urlcolor=black,
    pdfauthor={Fabián Belmar},
    pdftitle={Curriculum Vitae}
}

% Formato de secciones minimalista
\titleformat{\section}{\large\bfseries\scshape}{}{0em}{}[\vspace{2pt}\hrule\vspace{6pt}]
\titlespacing{\section}{0pt}{18pt}{8pt}

% Encabezado y pie
\pagestyle{fancy}
\fancyhf{}
\renewcommand{\headrulewidth}{0pt}
\fancyfoot[C]{\footnotesize \thepage}

% Comando para entradas
\newcommand{\cvitem}[2]{\noindent\textbf{#1}\hfill #2\\[2pt]}
\newcommand{\cvdetail}[1]{\noindent\hspace{0.3cm}#1\par\vspace{4pt}}

\begin{document}

% ============================================
% ENCABEZADO
% ============================================
\begin{center}
{\Large\textsc{Fabián Belmar}}\\[6pt]
{\small
\href{mailto:fbelmar@cepchile.cl}{fbelmar@cepchile.cl} ·
+56 9 76371608 ·
\href{https://belmarfabian.github.io}{belmarfabian.github.io}
}\\[3pt]
{\footnotesize
ORCID: \href{https://orcid.org/0000-0003-4239-1874}{0000-0003-4239-1874}
}
\end{center}

\vspace{12pt}

% ============================================
% POSICIÓN ACTUAL
% ============================================
\section{Posición Actual}

\cvitem{Centro de Estudios Públicos}{2025--}
\cvdetail{Investigador Asistente, Proyecto C22 - Métodos Digitales.}

% ============================================
% FORMACIÓN
% ============================================
\section{Formación Académica}

\cvitem{Universidad Adolfo Ibáñez}{2020--2025}
\cvdetail{Doctor en Procesos e Instituciones Políticas. Distinción máxima.}
\cvdetail{Tesis: ``Corrupción al margen del Estado. El caso de la FIFA''.}

\vspace{4pt}

\cvitem{Universidad de Talca}{2013--2018}
\cvdetail{Licenciado en Ciencia Política y Administración Pública. Distinción máxima.}

% ============================================
% EXPERIENCIA
% ============================================
\section{Experiencia en Investigación}

\cvitem{Instituto Gino Germani, U. de Buenos Aires}{2023--2024}
\cvdetail{Investigador pasante. Grupo de Estudios de Culturas Populares Contemporáneas.}

\vspace{4pt}

\cvitem{FONDECYT Regular \#1220004}{2022--2023}
\cvdetail{Asistente de investigación. IP: Dr. Mauricio Morales.}

\vspace{4pt}

\cvitem{Centro de Análisis Político, U. de Talca}{2018--2023}
\cvdetail{Coordinador de proyectos de investigación.}

\vspace{4pt}

\cvitem{FONDECYT Regular \#1180009}{2018--2020}
\cvdetail{Asistente de investigación. IP: Dr. Mauricio Morales.}

% ============================================
% PUBLICACIONES
% ============================================
\section{Publicaciones}

\subsection*{Artículos en revistas indexadas (WoS/Scopus)}

\begin{enumerate}[leftmargin=*, label={[\arabic*]}]
    \item Belmar, F., \& Mascareño, A. (2025). Corruption: An uneven field of research. \textit{Societies}, 15(7), 186.

    \item Belmar, F., \& Morales, M. (2025). Ethnic solidarity in the non-programmatic distribution of benefits. \textit{Ethnic and Racial Studies}.

    \item Belmar, F., Lara, B., \& Miguieles, J. (2024). Unpacking electoral competition. \textit{Policy Studies}.

    \item Belmar, F., Maldonado, F., \& Morales, M. (2024). In search of the missing link. \textit{Public Integrity}.

    \item Belmar, F., Morales, M., \& Villarroel, B. (2023). Writing constitution without parties? \textit{Politics}, 45(1).

    \item Belmar, F., Contreras, G., Morales, M., \& Troncoso, C. (2023). Demanding non-programmatic distribution. \textit{Policy Studies}.

    \item Belmar, F., \& Morales, M. (2022). Clientelism, turnout and incumbents' performance. \textit{Social Sciences}, 11(8), 361.

    \item Belmar, F., \& Morales, M. (2020). Clientelismo en la gestión local. \textit{Política y Sociedad}, 57(2), 567--591.
\end{enumerate}

\subsection*{Capítulos de libro}

\begin{enumerate}[leftmargin=*, label={[\arabic*]}]
    \item Belmar, F. (2024). Descentralización y participación ciudadana en Chile. En \textit{La Revolución Silente}.

    \item Belmar, F. (2020). Las tecnologías al servicio de la comunicación parlamentaria. En \textit{Parlamentos Conectados}. U. de Oviedo.

    \item Belmar, F. (2020). La estrategia de comunicación de los representantes. En \textit{Parlamentos Conectados}. U. de Oviedo.

    \item Belmar, F., Herrera, M., \& Obregón, M. (2018). El efecto de la elección de concejales sobre consejeros regionales. En \textit{Política Subnacional en Chile}. RIL Editores.
\end{enumerate}

% ============================================
% DOCENCIA
% ============================================
\section{Docencia}

\cvitem{Universidad de Talca}{2021--2022, 2025}
\cvdetail{Evaluación de Políticas; Ciencia de la Administración Pública; Métodos Cualitativos; Análisis de Datos.}

\vspace{4pt}

\cvitem{Universidad Diego Portales}{2025}
\cvdetail{Métodos Cuantitativos Introductorios; Métodos Cuantitativos Avanzados.}

\vspace{4pt}

\cvitem{Universidad Central}{2023--2024}
\cvdetail{Análisis de Datos Cuantitativos; Métodos Cuantitativos para la Investigación.}

\vspace{4pt}

\cvitem{Universidad Adolfo Ibáñez}{2022}
\cvdetail{Redes Sociales, Democracia y Ciudadanía Digital.}

% ============================================
% BECAS
% ============================================
\section{Becas y Distinciones}

\begin{itemize}[leftmargin=*, label={}]
    \item Beca Pasantía Internacional ANID, 2023--2024
    \item Beca Beneficios Complementarios ANID, 2021--2022
    \item Beca Doctorado Nacional ANID, 2020
    \item Mejor Egresado de la Generación, Universidad de Talca, 2018
\end{itemize}

% ============================================
% HABILIDADES
% ============================================
\section{Competencias}

\noindent\textbf{Software:} R, Python, SQL, STATA, SPSS, QGIS, ATLAS.ti, \LaTeX\\[2pt]
\noindent\textbf{Idiomas:} Español (nativo), Inglés (intermedio escrito)

% ============================================
% REFERENCIAS
% ============================================
\section{Referencias}

\noindent Mauricio Morales, Universidad de Talca. \href{mailto:mmoralesq@utalca.cl}{mmoralesq@utalca.cl}\\[2pt]
\noindent Aldo Mascareño, Universidad Adolfo Ibáñez. \href{mailto:amascareno@cepchile.cl}{amascareno@cepchile.cl}

\end{document}
